\section{The matrix the solver carries}
\label{sec:factorisation}

Proposition~\ref{prop:sub} gives the subproblem solution in closed form, and one
could implement the method by evaluating~\eqref{eq:subsol} afresh at every
iteration. That costs a $k \times k$ factorisation of the dual Hessian per step,
$O(nk^2 + k^3)$, and over the $O(n)$ steps a dual method takes it comes to
$O(n^4)$. The whole point of the Goldfarb--Idnani formulation is to replace it
with $O(n^2)$ per step, and it does so by carrying a single matrix through the
iteration and updating it.

\subsection{Two invariants}

Let $G = R_G\T R_G$ be the Cholesky factorisation with $R_G$ upper triangular,
and set $J_0 := R_G^{-1}$, again upper triangular. Then
$J_0 J_0\T = R_G^{-1} R_G^{-\top} = G^{-1}$ and $J_0\T G J_0 = I$. The solver's
state is a matrix $J \in \R^{n \times n}$ and an upper triangular
$R \in \R^{k \times k}$ satisfying
\begin{align}
  J\T G J &= I_n, \label{eq:inv1}\\
  J\T A &= \begin{bmatrix} R \\ 0 \end{bmatrix}. \label{eq:inv2}
\end{align}
Everything the iteration computes is a consequence of these two lines, so read
them slowly. Equation~\eqref{eq:inv1} is equivalent to
$J J\T = G^{-1}$ with $J$ nonsingular, and says that the columns of $J$ form a
basis of $\R^n$ orthonormal in the inner product $\langle v,w\rangle_G = v\T G w$.
Equation~\eqref{eq:inv2} says that in that basis the working-set normals are
triangular: partitioning
\begin{equation}
  J = \begin{bmatrix} J_1 & J_2 \end{bmatrix},
  \qquad J_1 \in \R^{n \times k},\quad J_2 \in \R^{n \times (n-k)},
  \label{eq:split}
\end{equation}
the lower block of~\eqref{eq:inv2} reads $J_2\T A = 0$.

At the cold start $\Active = \emptyset$, so~\eqref{eq:inv2} is vacuous and
$J = J_0$ serves. The updates of Section~\ref{sec:updates} maintain both
invariants by post-multiplying $J$ with orthogonal matrices, which preserves
\eqref{eq:inv1} exactly because $J\T G J = I$ is invariant under $J \mapsto JW$
for $W\T W = I$, while restoring the triangular shape of~\eqref{eq:inv2}.

\begin{proposition}[What the blocks mean]\label{prop:blocks}
Suppose \eqref{eq:inv1}--\eqref{eq:inv2} hold with $R$ nonsingular. Then
\begin{enumerate}
  \item[(i)] $R\T R = A\T G^{-1} A$, so $R$ is a Cholesky factor of the dual Hessian,
        unique up to the signs of its rows;
  \item[(ii)] $\range(J_2) = \nullsp(A\T)$, and the columns of $J_2$ are a
        $G$-orthonormal basis of it;
  \item[(iii)] $J_1 = G^{-1} A R^{-1}$ and $\range(J_1) = \range(G^{-1}A)$, the
        $G$-orthogonal complement of $\nullsp(A\T)$;
  \item[(iv)] $J_1 J_1\T = G^{-1}A(A\T G^{-1}A)^{-1}A\T G^{-1}$ and consequently
        \begin{equation}
          J_2 J_2\T = G^{-1} - G^{-1}A\bigl(A\T G^{-1}A\bigr)^{-1}A\T G^{-1}.
          \label{eq:reduced}
        \end{equation}
\end{enumerate}
\end{proposition}

\begin{proof}
(i) $A\T G^{-1} A = A\T J J\T A = R\T R$ by~\eqref{eq:inv2}, and two upper
triangular matrices with the same Gram matrix differ by a left factor
$\mathrm{diag}(\pm1)$.

(ii) $J_2\T A = 0$ places $\range(J_2) \subseteq \nullsp(A\T)$, with equality by
dimension since $J$ is nonsingular and $A$ has full column rank; the basis is
$G$-orthonormal because $J_2\T G J_2 = I_{n-k}$ is the trailing block
of~\eqref{eq:inv1}.

(iii) From~\eqref{eq:inv1}, $J^{-1} = J\T G$ and hence $J^{-\top} = G J$.
Inverting~\eqref{eq:inv2}, $A = J^{-\top}[R;0] = GJ[R;0] = G J_1 R$, so
$J_1 = G^{-1}AR^{-1}$, whose range is that of $G^{-1}A$ since $R$ is invertible.
That range is $G$-orthogonal to $\nullsp(A\T)$: for $v \in \nullsp(A\T)$,
$\langle G^{-1}Aw, v\rangle_G = w\T A\T v = 0$.

(iv) Substituting $J_1 = G^{-1}AR^{-1}$ and $R\T R = A\T G^{-1}A$,
\[
  J_1 J_1\T = G^{-1}A R^{-1} R^{-\top} A\T G^{-1}
            = G^{-1}A (R\T R)^{-1} A\T G^{-1}
            = G^{-1}A (A\T G^{-1}A)^{-1} A\T G^{-1},
\]
and \eqref{eq:reduced} follows from $J_1J_1\T + J_2J_2\T = JJ\T = G^{-1}$.
\end{proof}

The matrix in~\eqref{eq:reduced} is the \emph{reduced inverse Hessian} of the
working set: it inverts $G$ on the feasible subspace $\nullsp(A\T)$ and
annihilates the complement. Carrying $J$ therefore means carrying that object in
factored form, at no cost beyond the $n \times n$ array, which turns the search
directions of Section~\ref{sec:iteration} into single matrix--vector products.

\begin{remark}[Why fold the two factorisations together]\label{rem:fold}
An alternative state is $R_G$ together with an explicit orthogonal factor $Q$ of
the QR factorisation of $R_G^{-\top} A$. That holds the same information, $J$ being $J_0 Q$, and stores no more numbers. But every use of $J$ in
Section~\ref{sec:iteration} would become a triangular solve followed by an
orthogonal transformation, two operations where the folded form has one, and the
solve is the shape that does not reach a tuned kernel. The folded form has $J$
lose its triangularity after the first update, which is exactly the trade: $n^2$
storage and $2n^2$ flops per matrix--vector product, in exchange for the products
being $\texttt{gemv}$ at every iteration whatever the working set does.
\end{remark}

\subsection{The representation is not unique, and does not need to be}

Invariants~\eqref{eq:inv1}--\eqref{eq:inv2} do not determine $J$ and $R$. By
Proposition~\ref{prop:blocks}, $R$ is pinned only up to a left factor
$S = \mathrm{diag}(\pm1)$, whereupon $J_1 = G^{-1}AR^{-1}$ follows; and $J_2$ may
be any $G$-orthonormal basis of $\nullsp(A\T)$, so it is pinned only up to a right
factor $V$ with $V\T V = I_{n-k}$. The freedom is therefore exactly
\begin{equation}
  (J_1, J_2, R) \;\longmapsto\; (J_1 S,\; J_2 V,\; S R),
  \qquad S = \mathrm{diag}(\pm 1),\quad V\T V = I_{n-k}.
  \label{eq:freedom}
\end{equation}
Different orthogonal reductions realise different points of that family: a chain
of Givens rotations and a single Householder reflection annihilating the same
vector produce factors differing in sign, and a differently ordered sequence of
insertions and deletions produces a different $J_2$ outright. The following says
this does not matter, so the implementation may choose its reduction on grounds of
speed alone.

\begin{proposition}[Invariance of the iterates]\label{prop:invariance}
Let $(J,R)$ and $(\hat J,\hat R)$ both satisfy
\eqref{eq:inv1}--\eqref{eq:inv2} for the same $G$ and $A$. Then for every
$n_* \in \R^n$ the quantities
\begin{equation}
  d = J\T n_*, \qquad z = J_2 d_2, \qquad r = R^{-1} d_1
  \label{eq:dzr}
\end{equation}
with $d = (d_1,d_2)$ split at $k$, satisfy $\hat z = z$ and $\hat r = r$.
Moreover, if the two representations are related by~\eqref{eq:freedom} with
$V = I$, the equalities hold exactly in IEEE~754 arithmetic.
\end{proposition}

\begin{proof}
By Proposition~\ref{prop:blocks}(i) and $A\T J = [R\T\ 0]$ we have
$A\T G^{-1} n_* = A\T J J\T n_* = R\T d_1$, so
\begin{equation}
  r = R^{-1}d_1 = R^{-1}R^{-\top}R\T d_1 = (R\T R)^{-1} A\T G^{-1} n_*
    = \bigl(A\T G^{-1}A\bigr)^{-1} A\T G^{-1} n_* ,
  \label{eq:rclosed}
\end{equation}
which mentions only $A$, $G$ and $n_*$. Likewise $z = J_2 J_2\T n_*$, since
$d_2 = J_2\T n_*$, and by~\eqref{eq:reduced}
\begin{equation}
  z = \Bigl(G^{-1} - G^{-1}A\bigl(A\T G^{-1}A\bigr)^{-1}A\T G^{-1}\Bigr) n_*,
  \label{eq:zclosed}
\end{equation}
again free of the representation. For the last claim, $\hat R = SR$ gives
$\hat d_1 = S d_1$ and $\hat r = (SR)^{-1}Sd_1 = R^{-1}S^{-1}Sd_1$, in which the
only floating-point operations beyond those forming $r$ are multiplications by
$\pm1$; these are exact, so the computed $\hat r$ and $r$ agree bit for bit. The
same argument applies to $\hat z = \sum_{j>k}(s_j J_{\cdot j})(s_j d_j)$.
\end{proof}

\begin{remark}[What the proposition does and does not settle]
It settles that the substitution is exact in the sense that matters: the
iteration's decisions (which constraint is most violated, whether the step is full
or partial, which constraint is dropped) are functions of $z$ and $r$
alone, and those are functions of $(G,A,n_*)$ alone. It does not settle that two
implementations follow the same trajectory in finite precision, because the
rounding \emph{within} the reductions differs and the decisions are
discontinuous functions of the computed values. That is an empirical question,
and settling it properly needs two independent implementations to compare
trajectories against each other. What is checkable against one, and is checked in
the Reproducibility note, is the weaker statement: that the $z$ and $r$ the solver
computes agree with the representation-free closed forms~\eqref{eq:zclosed}
and~\eqref{eq:rclosed} to machine precision.
\end{remark}
